% Auto-generated by extract_cwe_cve_examples.py
% Brief version for thesis writing.

\paragraph{CWE-863: Incorrect Authorization（授权不当）}
授权逻辑存在错误或缺失，导致未授权主体获得本应受限的资源访问。\
\noindent\textbf{CVE-2019-11247}（Kubernetes）：The Kubernetes kube-apiserver mistakenly allows access to a cluster-scoped custom resource if the request is made as if the resource were namespaced. Authorizations for the resource accessed in this manner are enforced using roles and role bindings within the ...（成因映射：安全逻辑缺陷/2.1 访问控制缺陷）\

\paragraph{CWE-269: Improper Access Control (Generic)（权限/访问控制管理不当）}
权限管理或访问控制配置不当，使得主体获得超过最小权限原则的能力。\
\noindent\textbf{CVE-2022-29179}（Cilium）：Cilium is open source software for providing and securing network connectivity and loadbalancing between application workloads. Prior to versions 1.9.16, 1.10.11, and 1.11.15, if an attacker is able to perform a container escape of a container running as root ...（成因映射：配置错误/1.1 权限配置过宽）\

\paragraph{CWE-22: Improper Limitation of a Pathname to a Restricted Directory（路径遍历）}
对路径/目录边界限制不足，攻击者可利用相对路径等手段访问受限目录之外的文件。\
\noindent\textbf{CVE-2022-24730}（Argo CD）：Argo CD is a declarative, GitOps continuous delivery tool for Kubernetes. Argo CD starting with version 1.3.0 but before versions 2.1.11, 2.2.6, and 2.3.0 is vulnerable to a path traversal bug, compounded by an improper access control bug, allowing a malicious...（成因映射：代码注入/4.2 路径遍历）\

\paragraph{CWE-287: Improper Authentication（认证不当）}
认证流程不完善（例如绕过、校验缺失或错误），导致身份校验被规避。\
\noindent\textbf{CVE-2022-29165}（Argo CD）：Argo CD is a declarative, GitOps continuous delivery tool for Kubernetes. A critical vulnerability has been discovered in Argo CD starting with version 1.4.0 and prior to versions 2.1.15, 2.2.9, and 2.3.4 which would allow unauthenticated users to impersonate ...（成因映射：安全逻辑缺陷/2.2 身份验证绕过）\

